%% THERMODYNAMIQUE %%
%
\section{Temp\'erature potentielle \'equivalente }

\itemstyle{ Op\'erateur~: } {\sf TPotEq}
\index{TPotEq}

%%%%%OBJET%%%%%
\itemstyle{ Objet~: }
Calcul de la temp\'erature potentielle \'equivalente de l'air.
La temp\'erature \'equivalente est la temp\'erature que la parcelle
d'air humide atteindrait si on condensait toute la vapeur
d'eau, \`a pression constante. On se ram\`ene ensuite \`a 1000 hPa
selon une transformation adiabatique, pour obtenir une temp\'erature
potentielle.

%%%%%ARGUMENTS%%%%%
\bigskip
%%%%%ENTREES%%%%
\begin{tabarg}
\multicolumn{3}{l}{\titretab{Entr\'ees}}\\
\hline
$T_s$ & Vecteur, {\tt float} & Temp\'erature statique ($K$ ou $^\circ C$) \\
$\theta $ & Vecteur, {\tt float} & Temp\'erature potentielle ($K$ ou $^\circ C$) \\
$r$ & Vecteur, {\tt float} & Rapport de m\'elange en eau liquide ($g/kg$)\\
$c_{pa}$ & Coeff., {\tt float} & Chaleur massique \`a pression constante de l'air sec\\
\end{tabarg}

\medskip
%%%%%SORTIES%%%%
\begin{tabarg}
\multicolumn{3}{l}{\titretab{Sorties}}\\
\hline
$\theta_e$ & Vecteur, {\tt float} & Temp\'erature potentielle \'equivalente (m\^eme unit\'e que $T_s$)\\
\end{tabarg}


%%%%%FORMULE%%%%
\itemstyle{ Formule }


% calcul de la temp\'erature potentielle equivalente 

\[
\theta_e = \theta .\left( 1 + r. \frac{L}{c_{pa}.T_s} \right) 
\]
avec $ L = 3136.17 - 2.34\times T_s $ (pour $T_s$ exprim\'e en $K$).\\


%%%%% COMMENTAIRE %%%

%\itemstyle{ Commentaire~: }


%%%%%AUTEUR%%%%%
\itemstyle{ Auteur~:}
CNRM/GMEI/TRAMM

%%%%%REFERENCES%%%%%
\itemstyle{ R\'ef\'erences~:}
Directement copi\'ee de la routine du CAM, qui est identique 
\`a l'algorithme que P. Durand cite dans le recueil de formules
\'etabli \`a l'occasion de PYREX.
